\section{Discussion and Limitations}
\label{sec:discussion}

\subsection{Physical Implications for Metamaterial Design}
The results presented in Sections~\ref{sec:results} and \ref{sec:steering} establish that an anisotropic characteristic length tensor offers an effective mechanism for tuning directional wave propagation. In classical phononic crystals, directional filtering is typically achieved through geometric asymmetry of inclusions or complex multi-phase material architectures. Here, we demonstrate that in materials possessing intrinsic microstructural texture---such as cold-rolled alloys, oriented fibrous media, or architected sub-lattices---directional filtering can be manipulated solely by adjusting the orientation of the microstructural axes relative to the lattice coordinate frame.

\subsection{Distinction Between Directional Stop Bands and Complete Band Gaps}
A critical finding of this study is the strict adherence to the gap hierarchy $\Delta[\mathrm{leg}] \ge \Delta[\mathrm{path}] \ge \Delta[\mathrm{complete}]$. In the homogeneous Case H configuration, while substantial directional stop bands open along specific high-symmetry paths (e.g., $\Delta_{GX} = 0.0896$ at $\mathrm{AR} = 10, \theta = 30^\circ$, and $\Delta_{GX} = 0.0431$ at $\theta = 45^\circ$), the global complete band gap remains strictly non-positive ($\Delta_{\mathrm{complete}} \le -0.3758$). Propagating modes in other directions (specifically along $X$--$M$ and $M$--$\Gamma$) span the stop-band frequencies. Omnidirectional complete band gaps require periodic impedance contrast, such as the two-phase composite inclusions planned in Case C.

\subsection{Explicit Methodological Scope and Limitations}
To ensure complete scientific transparency, we explicitly enumerate the methodological boundaries and open items associated with this limited-scope manuscript:
\begin{enumerate}
  \item \textbf{Status of External Validation Benchmarks (Gate G3 / PCR1):}
  Formal quantitative validation against the 1D published transfer-matrix benchmarks of Li et al.\ (2023, 2024) \cite{li2023anchorA,li2024anchorB} and Mishra et al.\ (2026) \cite{mishra2026anchorC} is currently designated as BLOCKED under Gate G3. While qualitative agreement and internal consistency are fully established, discrepancies in external transfer-matrix normalizations and state-vector formulations are undergoing technical remediation. Consequently, Benchmarks B1, B2, and B3 remain unvalidated, and Benchmark B6 is classified as PARTIAL.
  \item \textbf{Status of Case C Phononic Crystal Study (Study S2):}
  Composite phononic inclusions with material contrast across the unit cell are currently BLOCKED pending resolution of interface continuity conditions for $C^1$ elements across bimaterial discontinuities.
  \item \textbf{Open Technical Variations:}
  Technical variations TV1 (normal derivative Bloch coupling), TV6 (interface gradient traction jump conditions), TV12 (shear-locking in thin aspect ratios), TV14 (high-frequency spurious optical mode filtering), and TV18 (three-dimensional generalization) remain open research questions.
  \item \textbf{Kinematic and Constitutive Assumptions:}
  The analysis is restricted to linear, isothermal, non-piezoelectric, two-dimensional plane-strain conditions under infinitesimal deformations. Non-linear strain-gradient effects, thermal dissipation, and 3D out-of-plane polarization couplings are beyond the scope of this investigation.
\end{enumerate}
